\documentclass[12pt]{letter}
\usepackage{geometry}
\geometry{margin=1in}
\usepackage{hyperref}

\begin{document}

\begin{letter}{%
Editor-in-Chief\\
\emph{IEEE Transactions on Computational Social Systems}}

\opening{Dear Editor,}

I am pleased to submit the manuscript ``Geometric Prediction of Economic Behavior: Cross-Domain Validation Across Game Theory and Prospect Theory'' for consideration in \emph{IEEE Transactions on Computational Social Systems}.

The paper demonstrates that a geometric framework---modeling economic decisions as pathfinding on a nine-dimensional manifold with Mahalanobis distance as the cost metric---achieves quantitative prediction of human economic behavior across two traditionally separate literatures: strategic game theory and individual prospect theory.  A single three-parameter covariance matrix, calibrated on bargaining game data via automated structural fuzzing, predicts all 16 heterogeneous targets within tolerance (MAE = 2.70\%, 16/16 pass), including six prospect theory targets that are genuine out-of-sample predictions.

The paper makes five principal contributions:
\begin{enumerate}
    \item \textbf{Cross-domain unification:} A single calibrated metric predicts both strategic interaction (ultimatum, dictator, public goods games) and individual choice under risk (Kahneman--Tversky prospect theory problems), with no domain-specific recalibration.
    \item \textbf{16/16 prediction targets at 2.70\% MAE:} The model passes all targets with a degrees-of-freedom ratio of 5.3:1 (16 targets, 3 parameters), exceeding the standard threshold for genuine prediction.
    \item \textbf{Money-zero finding:} The monetary dimension ($d_1$) has zero sensitivity---removing it does not change predictions.  Economic behavior is driven by social, identity, and epistemic dimensions, not money.
    \item \textbf{Systematic baseline comparison:} Direct comparison against Cumulative Prospect Theory (6 parameters), Fehr--Schmidt inequity aversion (2+ parameters), and expected-value maximization (0 parameters) on a shared target set, showing the geometric model achieves higher accuracy with fewer parameters and broader coverage.
    \item \textbf{Cost-dependent temperature:} A principled mechanism linking population-level choice heterogeneity to the geometric cost gap between options, resolving the Allais paradox within the framework.
\end{enumerate}

This work is relevant to \emph{IEEE TCSS} because it provides a computational framework for predicting social behavior that unifies game-theoretic and decision-theoretic phenomena, validated against published experimental data from multiple independent research groups.

The manuscript has not been published previously and is not under consideration elsewhere.  No human participants were involved; this is a computational modeling paper using published experimental data.  The software implementation is publicly available under the MIT open-source license.

Generative AI tools (Claude, Anthropic) were used for manuscript preparation and formatting; all intellectual content, mathematical framework, and experimental design are the author's own.

Thank you for considering this submission.

\closing{Sincerely,\\[2em]
Andrew H. Bond\\
Senior Lecturer, Department of Computer Engineering\\
San Jos\'{e} State University\\
\href{mailto:andrew.bond@sjsu.edu}{andrew.bond@sjsu.edu}\\
ORCID: 0009-0003-2599-6158}

\end{letter}
\end{document}
